\section{Methodology}

\subsection{Overall Architecture}

In this section, we present the overall architecture of the proposed Unified Dual-Mask Physical Model for non-Lambertian light field depth estimation. Different from previous works [1, 2] that treat all surface reflectances uniformly or rely on heuristic post-processing, our framework explicitly integrates a physical reflectance model into a dual-path neural architecture. As illustrated in Fig. 1, the network comprises a Three-Layer Angular Signal Decomposition \& Mask Generation Module, a Lambertian EPI Branch, a Non-Lambertian Geometric Branch, and a Balanced Projection \& Mask-Weighted Fusion Module. This design enables the adaptive routing and specialized processing of features based on the underlying material properties, ensuring robust depth recovery across heterogeneous scenes.

The input to our model consists of densely sampled light field data with a $9 \times 9$ angular resolution, represented as a tensor $\mathcal{I} \in \mathbb{R}^{B \times 9 \times 9 \times H \times W \times 3}$, where $B$ is the batch size, and $H$ and $W$ denote the spatial height and width, respectively. To capture comprehensive epipolar geometry, we extract 4-directional Epipolar Plane Images (EPIs) from the angular grid, serving as the foundational input for the Lambertian EPI Branch. Prior to feeding the data into the network, we apply tone mapping augmentation to simulate diverse illumination conditions and enhance the generalization to varying exposure levels. This preprocessing step ensures that the subsequent physical decomposition remains stable under complex real-world lighting variations. The preprocessed light field tensor is then simultaneously directed to the decomposition module and the Lambertian branch for parallel feature extraction.

The core data flow follows a "decomposition $\rightarrow$ dual-branch parallel extraction $\rightarrow$ mask-guided fusion" topology. Initially, the light field data enters the Three-Layer Angular Signal Decomposition \& Mask Generation Module. Motivated by the physical image formation process, this module decomposes the angular signal into three components: the direct current (DC) term representing ambient illumination, the linear term encoding parallax, and the residual term capturing material-specific reflectance. Mathematically, the angular intensity $I(u,v)$ at viewpoint $(u,v)$ is modeled as $I(u,v) = I_{DC} + a \cdot u + b \cdot v + \varepsilon(u,v)$, where $I_{DC}$ is the ambient light, $a$ and $b$ are linear coefficients corresponding to depth, and $\varepsilon(u,v)$ is the residual reflecting material properties. From this decomposition, we compute geometric angular features, such as the angular gradient $P(x,y) = \sqrt{a^2+b^2}$ and correlation coherence, yielding Geometric Angular Features $\mathcal{F}_{geo} \in \mathbb{R}^{B \times H \times W \times C_{geo}}$. These features are fed into a classifier to generate a Pixel-wise Dual-Mask $\mathcal{M} \in \mathbb{R}^{B \times 2 \times H \times W}$, which partitions the scene into Lambertian and non-Lambertian regions. Guided by $\mathcal{M}$, the data flows into two specialized branches. The Lambertian EPI Branch leverages the linear EPI structures and incorporates continuous depth priors to extract Lam Depth Features $\mathcal{F}_{lam} \in \mathbb{R}^{B \times C_{lam} \times H \times W}$, enhancing surface smoothness. Conversely, the Non-Lambertian Geometric Branch processes $\mathcal{F}_{geo}$ to capture complex X-shaped EPI patterns without relying on the violated linear assumption, producing NL Depth Features $\mathcal{F}_{nl} \in \mathbb{R}^{B \times C_{nl} \times H \times W}$.

To integrate the heterogeneous features from the dual branches, we employ the Balanced Projection \& Mask-Weighted Fusion Module. Different from previous works [3, 4] that suffer from severe channel imbalance and feature confusion when directly concatenating Lambertian and non-Lambertian representations, our module explicitly aligns the feature spaces before aggregation. Specifically, we project $\mathcal{F}_{lam}$ and $\mathcal{F}_{nl}$ into a unified balanced dimension ($C_{fuse} = 96$) using $1 \times 1$ convolutions. The projected features are then aggregated via pixel-wise weighting guided by the Pixel-wise Dual-Mask $\mathcal{M}$:
\begin{equation}
\label{eq:eq1}
\mathcal{F}_{fuse} = \mathcal{M}_{lam} \odot \text{Proj}_{lam}(\mathcal{F}_{lam}) + \mathcal{M}_{nl} \odot \text{Proj}_{nl}(\mathcal{F}_{nl})
\end{equation}
where $\odot$ denotes the Hadamard product, $\mathcal{M}_{lam}$ and $\mathcal{M}_{nl}$ are the respective channels of the dual-mask, and $\text{Proj}(\cdot)$ represents the projection operation. The fused depth feature $\mathcal{F}_{fuse} \in \mathbb{R}^{B \times 96 \times H \times W}$ is subsequently passed through a Depth Decoder, which consists of progressive upsampling and convolutional layers, to regress the Final Depth Map $\mathcal{D} \in \mathbb{R}^{B \times 1 \times H \times W}$. This mask-weighted fusion mechanism ensures that the final depth prediction strictly adheres to the physical characteristics of each local region, yielding highly accurate and continuous depth maps for downstream 3D video systems.

\subsection{Three-Layer Angular Signal Decomposition and Mask Generation Module}

Conventional light field depth estimation methods [5, 6] predominantly rely on the Lambertian assumption, treating all surface reflectances uniformly. This simplification leads to severe depth distortions in regions with specular reflections or complex materials. To address this limitation, we develop the Three-Layer Angular Signal Decomposition and Mask Generation Module. Unlike previous approaches [7, 8] that implicitly learn material differences or apply heuristic post-processing, this module explicitly integrates a physical reflectance model to decouple the angular signals. The primary objective is to decompose the input light field into distinct physical components, extract geometry-aware angular features, and generate pixel-level dual masks. These masks serve as routing mechanisms for the subsequent dual-path architecture, ensuring that Lambertian and non-Lambertian regions are processed by their respective specialized branches.

The module takes the densely sampled $9 \times 9$ light field data $\mathcal{I} \in \mathbb{R}^{B \times 9 \times 9 \times H \times W \times 3}$ as input. For a specific spatial location $(x,y)$, the angular signal across the view grid $(u,v) \in \{-4, \dots, 4\}^2$ is denoted as $I_{x,y}(u,v)$. Based on the physical image formation model, the angular intensity variation can be approximated by a first-order Taylor expansion combined with a residual term. We formulate the three-layer decomposition as:
\begin{equation}
\label{eq:eq2}
I_{x,y}(u,v) = I_{dc}(x,y) + a(x,y) \cdot u + b(x,y) \cdot v + \varepsilon_{x,y}(u,v)
\end{equation}
where $I_{dc}(x,y)$ represents the direct current (DC) component, corresponding to the ambient illumination and base albedo. The linear terms $a(x,y)$ and $b(x,y)$ denote the angular gradients along the $u$ and $v$ axes, which are intrinsically linked to the local disparity and depth. The residual term $\varepsilon_{x,y}(u,v)$ captures the high-frequency angular variations caused by non-Lambertian reflections, such as specular highlights, as well as sensor noise. We specifically choose the residual term $\varepsilon_{x,y}(u,v)$ to reflect material attributes because non-Lambertian phenomena violate the linear intensity variation assumption along the epipolar line; thus, the energy of this residual serves as a direct physical indicator of surface reflectance complexity. To solve for these components, we employ a least-squares fitting approach over the $9 \times 9$ angular grid for each spatial pixel, effectively separating the diffuse and specular reflectance properties.

Following the decomposition, we compute geometric angular features to quantify the material properties at each pixel. The magnitude of the linear angular gradient is calculated as:
\begin{equation}
\label{eq:eq3}
P(x,y) = \sqrt{a(x,y)^2 + b(x,y)^2}
\end{equation}
A higher value of $P(x,y)$ indicates a strong linear epipolar structure, characteristic of Lambertian surfaces. Conversely, non-Lambertian regions exhibit significant energy in the residual term. We define the angular gradient feature, denoted as $F_{ang}(x,y)$, by combining the gradient magnitude with the variance of the residual term $\text{Var}(\varepsilon_{x,y})$ and the angular correlation coherence. Additionally, the coefficient of variation is utilized to normalize the residual energy against the base intensity, ensuring robustness to varying illumination conditions. These computed features are then concatenated and passed through a multi-layer perceptron to generate the final pixel-wise dual masks, completing the physical decomposition and routing preparation. With the complete pipeline established, we now proceed to empirically validate the proposed framework and detail the experimental setup used for performance assessment.